\begin{pastexam}[label={px:df3-100-ds-02}]{地方三等資訊處理 100 年 第二題}
（申論；20 分）

假設有一整數資料陣列 B[0..7]，裡面儲存 8 個整數數值分別為\{25, 57, 86, 37, 12, 92, 48, 33\}。今欲對此陣列進行
由小到大排序：
\begin{itemize}[leftmargin=2.2em, nosep]
\item[㈠] 試寫出氣泡浮昇排序（bubble sort）演算法或函式。（10 分）
\item[㈡] 將排序過程中每一回合（iteration）陣列內容的變化情形寫出。（10 分）
\end{itemize}
\end{pastexam}

\begin{pastsoln}
答案：㈡ 五個回合後排好，第六回合沒有交換而結束；各回合結果見下表。

\textbf{㈠} 每一回合從左到右比較相鄰兩數，左邊比較大就交換，這一回合結束時最大的數被換到最右邊；下一回
合少看最後一格。 某一回合完全沒有交換，表示已經排好，提早結束（資結書第~14.2~節、虛擬碼~14.1）。 依
題目用 $0$ 起算的陣列寫成 C 函式：
\begin{minted}[fontsize=\small]{c}
void bubble_sort(int B[], int n)
{
    for (int i = 0; i < n - 1; i++) {        /* 第 i+1 回合 */
        int swapped = 0;
        for (int j = 0; j < n - 1 - i; j++) {
            if (B[j] > B[j + 1]) {           /* 相鄰兩數順序不對就交換 */
                int t = B[j]; B[j] = B[j + 1]; B[j + 1] = t;
                swapped = 1;
            }
        }
        if (!swapped)                        /* 這回合沒有交換：已排好 */
            break;
    }
}
\end{minted}
第 $i + 1$ 回合比較 $n - 1 - i$ 次，結束時 $B[n-1-i..n-1]$ 是最大的 $i + 1$ 個數，已在最後位置。 最壞時間
$O(n^2)$，已排好時只要一回合 $O(n)$。

\textbf{㈡} 以第一回合為例逐步看（粗體是被換到右邊的數）：
\begin{itemize}
\item $25, 57$ 不換；$57, 86$ 不換；
\item $86 > 37$ 換：$25, 57, 37, \mathbf{86}, 12, 92, 48, 33$；
\item $86 > 12$ 換：$25, 57, 37, 12, \mathbf{86}, 92, 48, 33$；
\item $86, 92$ 不換；
\item $92 > 48$ 換：$25, 57, 37, 12, 86, 48, \mathbf{92}, 33$；
\item $92 > 33$ 換：$25, 57, 37, 12, 86, 48, 33, \mathbf{92}$。
\end{itemize}
每一回合結束時的陣列（直線右邊是已經就定位的數）：
\begin{center}
\begin{tabular}{lll}
\toprule
回合 & 陣列 $B[0..7]$ & 交換次數 \\
\midrule
開始 & $25, 57, 86, 37, 12, 92, 48, 33$ & \\
1 & $25, 57, 37, 12, 86, 48, 33 \mid 92$ & 4 \\
2 & $25, 37, 12, 57, 48, 33 \mid 86, 92$ & 4 \\
3 & $25, 12, 37, 48, 33 \mid 57, 86, 92$ & 3 \\
4 & $12, 25, 37, 33 \mid 48, 57, 86, 92$ & 2 \\
5 & $12, 25, 33 \mid 37, 48, 57, 86, 92$ & 1 \\
6 & $12, 25 \mid 33, 37, 48, 57, 86, 92$ & 0（提早結束） \\
\bottomrule
\end{tabular}
\end{center}
共交換 $4 + 4 + 3 + 2 + 1 = 14$ 次，等於原陣列的逆序對個數。 若不寫提早結束，還會再做第 $7$ 回合
（只比較 $12, 25$，不換），陣列不變。 最後結果 $12, 25, 33, 37, 48, 57, 86, 92$。
\end{pastsoln}
